\tikzset{>=latex} % for LaTeX arrow head

\colorlet{myred}{red!65!black}
\colorlet{mygreen}{green!60!black}
\colorlet{mydarkred}{red!50!black}
\colorlet{mydarkblue}{blue!40!black}
\tikzstyle{measure}=[{Latex[length=4,width=3]}-{Latex[length=4,width=3]},line width=0.4,mydarkblue]
\tikzstyle{ground}=[preaction={fill,top color=black!10,bottom color=black!5,shading angle=20},
                    fill,pattern=north east lines,draw=none,minimum width=0.3,minimum height=0.6]
\tikzstyle{mass}=[line width=0.6,red!30!black,fill=red!40!black!10,rounded corners=1,
                  top color=red!40!black!20,bottom color=red!40!black!10,shading angle=20]
\tikzstyle{rope}=[brown!70!black,line width=1.2,line cap=round] %very thick
\tikzstyle{force}=[->,myred,thick,line cap=round]
\tikzstyle{unit}=[->,mygreen,thick,line cap=round]
\tikzset{fieldlines/.style={mydarkred,decoration={markings,mark=at position #1 with {\arrow{latex}}},
                            postaction={decorate},line width=0.7},
         fieldlines/.default=0.55}
\newcommand{\vbF}{\vb{F}}

% GRAVITATIONAL ATTRACTION


\begin{tikzpicture}
  \def\d{4.1}  % distance r
  \def\r{0.3}  % small radius ball
  \def\R{0.38} % large radius ball
  \def\F{0.95} % force magnitude
  \def\u{0.4}  % unit vector magnitude
  \def\ang{5}  % angle lin
  \coordinate (M1) at (0,0);
  \coordinate (M2) at (\ang:\d);
  
  % SETUP
  \draw[measure,-]
    (M1) --++ (\ang+90:1.6*\R)
    (M2) --++ (\ang+90:1.6*\R);
  \draw[measure] ($(M1)+(\ang+90:1.35*\R)$) --++ (M2)
    node[midway,circle,fill=white,inner sep=0] {$r$};
  \draw[dashed] (M1) -- (M2);
  \draw[mass] (M1) circle(\r) node[scale=0.9] {$m_1$};
  \draw[mass] (M2) circle(\R) node[scale=0.9] {$m_2$};
  
  % FORCES
  \draw[force] (M1)++(\ang-8:0.9*\r) --++ (\ang:\F)
    node[anchor=90,inner sep=4] {$\vbF_{21}$}; % = -G\dfrac{m_1m_2}{r^2}\vu{r}_{21}
  \draw[force] (M2)++(\ang-172:0.9*\R) --++ (\ang:-\F)
    node[anchor=130,inner sep=3] {$\vbF_{12}$}; % = -G\dfrac{m_1m_2}{r^2}\vu{r}_{12}
  \draw[unit] (M1)++(\ang+8:\r+0.02) --++ (\ang:\u)
    node[pos=0.8,above=-1] {$\vu{r}_{12}$};
  \draw[unit] (M2)++(\ang+172:\R+0.02) --++ (\ang:-\u)
    node[pos=0.6,above=-1] {$\vu{r}_{21}$};
  
\end{tikzpicture}
